\section{Deterministic selection and positivity}\label{sec:selection}

We now extract a single deterministic approximation from the random model.
The two required high-probability properties are a partition-function upper
bound and asymptotic freeness.

\begin{lemma}[Exposure bound]\label{lem:exposure}
Let $w\in\F_2\setminus\{e\}$ be a reduced word of length $\ell$.  For the
homomorphism generated by two independent uniform permutations of $[n]$,
there is a constant $C_w$ such that
\[
 \mathbb E|\Fix(\sigma_n(w))|\le C_w
\]
for all sufficiently large $n$.  Hence
$|\Fix(\sigma_n(w))|/n\to0$ in probability.
\end{lemma}

\begin{proof}
Write $w=t_\ell\cdots t_1$ with
$t_i\in\{a,a^{-1},b,b^{-1}\}$.  Starting at a fixed $v=x_0$, expose the path
$x_i=\sigma_n(t_i)x_{i-1}$.  As long as $x_0,\ldots,x_{i-1}$ are distinct,
reducedness prevents step $i$ from being the forced reverse of the preceding
edge.  At most $i-1$ values of the relevant permutation have been exposed,
so $x_i$ is uniform among at least $n-\ell$ available vertices.  At most $i$
of them cause the first collision.  Thus, for $n>\ell$,
\[
 \mathbb P(\sigma_n(w)v=v)
 \le\sum_{i=1}^\ell\frac{i}{n-\ell}
 =\frac{\ell(\ell+1)}{2(n-\ell)}.
\]
For $n\ge2\ell$ this is at most $\ell(\ell+1)/n$.  Summing over $v$ gives the
claim with $C_w=\ell(\ell+1)$.  Markov's inequality gives convergence in
probability after division by $n$.
\end{proof}

\begin{proposition}[Deterministic uniform approximation]\label{prop:selection}
There is a deterministic homomorphic sofic approximation
$\Sigma_{\unif}=(\tau_m:\F_2\to\Sym(V_m))$ such that
\begin{equation}\label{eq:upper-selected}
 \limsup_{m\to\infty}\frac1{|V_m|}
 \log Z_{G_{\tau_m}}(3)\le C_3.
\end{equation}
\end{proposition}

\begin{proof}
By \Cref{prop:annealed}, for every fixed $\varepsilon>0$ and all sufficiently
large $n$,
\[
 \mathbb E Z_{G_n}(3)\le\exp(n(C_3+\varepsilon/2)).
\]
Markov's inequality therefore gives
\begin{equation}\label{eq:markov-Z}
 \mathbb P\left(Z_{G_n}(3)>\exp(n(C_3+\varepsilon))\right)
 \le e^{-n\varepsilon/2}.
\end{equation}

Enumerate the nonidentity elements as $w_1,w_2,\ldots$.  At stage $m$, set
$\varepsilon=1/m$.  By \eqref{eq:markov-Z}, the desired partition bound has
probability tending to one as $n\to\infty$.  By \Cref{lem:exposure} and a
union bound over $w_1,\ldots,w_m$, the event
\[
 |\Fix(\sigma_n(w_j))|<n/m,
 \qquad 1\le j\le m,
\]
also has probability tending to one.  Choose a strictly increasing $n_m$ for
which the intersection of these events is nonempty, choose one realization,
and call it $\tau_m$.

Each $\tau_m$ is a homomorphism.  For fixed $w_j\ne e$, the selected
fixed-point fraction is at most $1/m$ for all $m\ge j$, so it tends to zero.
Thus $(\tau_m)$ is a homomorphic sofic approximation.  Its selected
partition bounds give \eqref{eq:upper-selected}.
\end{proof}

Loops must be removed before applying a bounded-degree independence estimate.
This costs only a vanishing fraction.

\begin{proposition}[Positive lower bound]\label{prop:positive}
The approximation $\Sigma_{\unif}$ from \Cref{prop:selection} satisfies
\[
 \liminf_{m\to\infty}\frac1{|V_m|}
 \log Z_{G_{\tau_m}}(3)\ge\frac15\log4.
\]
Consequently,
\[
 \frac15\log4
 \le h_{\Sigma_{\unif}}(X_3)\le C_3.
\]
\end{proposition}

\begin{proof}
Let $L_m$ be the looped vertices of $G_{\tau_m}$.  A loop can only occur at a
fixed point of $\tau_m(a)$ or $\tau_m(b)$, so soficity gives
\[
 |L_m|
 \le |\Fix(\tau_m(a))|+|\Fix(\tau_m(b))|
 =o(|V_m|).
\]
Delete $L_m$ and replace parallel edges by single edges.  The resulting
simple graph has maximum degree at most four: each vertex has at most an
incoming and an outgoing neighbour for each of the two generator
permutations.  The greedy algorithm in a graph of maximum degree $\Delta$
selects an independent set of size at least the number of vertices divided by
$\Delta+1$.  Hence $G_{\tau_m}$ has an independent set $I_m$ with
\[
 |I_m|\ge\frac{|V_m|-|L_m|}{5}.
\]

For every vertex of $I_m$, choose independently either zero or one of the
three active colours, and fill $V_m\setminus I_m$ by zero.  These are
$4^{|I_m|}$ distinct exact models.  Therefore
\[
 \liminf_m\frac1{|V_m|}\log Z_{G_{\tau_m}}(3)
 \ge\liminf_m\frac{|I_m|}{|V_m|}\log4
 \ge\frac15\log4.
\]
The entropy bounds follow from \Cref{cor:model-formula,prop:selection}.
\end{proof}

\begin{proof}[Proof of \Cref{thm:main}]
The bipartite lower bound is \Cref{prop:bip-lower}.  The uniform upper and
positive lower bounds are \Cref{prop:selection,prop:positive} together with
\Cref{cor:model-formula}.  The strict middle inequality is
\Cref{lem:strict-gap}.  Combining them gives \eqref{eq:main-chain}.
\end{proof}

